\documentclass[11pt,letterpaper]{article}

\usepackage[utf8]{inputenc}
\usepackage[T1]{fontenc}
\usepackage[spanish]{babel}
\usepackage[margin=2.5cm]{geometry}
\usepackage{graphicx}
\usepackage{xcolor}
\usepackage{listings}
\usepackage{textcomp}
\usepackage{url}
\usepackage[hidelinks]{hyperref}



\definecolor{fondo}{RGB}{246,248,250}
\definecolor{clave}{RGB}{0,92,197}
\definecolor{texto}{RGB}{181,80,0}
\definecolor{comentario}{RGB}{106,115,125}

\lstset{
  language=SQL,
  upquote=true,
  basicstyle=\ttfamily\small,
  keywordstyle=\color{clave}\bfseries,
  stringstyle=\color{texto},
  commentstyle=\color{comentario}\itshape,
  backgroundcolor=\color{fondo},
  frame=single,
  rulecolor=\color{black!15},
  breaklines=true,
  columns=fullflexible,
  keepspaces=true,
  showstringspaces=false,
  morekeywords={ENUM,AUTO_INCREMENT,LIMIT,DECIMAL,DATE,IF,EXISTS},
  literate={á}{{\'a}}1 {é}{{\'e}}1 {í}{{\'i}}1 {ó}{{\'o}}1 {ú}{{\'u}}1
           {Á}{{\'A}}1 {É}{{\'E}}1 {Í}{{\'I}}1 {Ó}{{\'O}}1 {Ú}{{\'U}}1
           {ñ}{{\~n}}1 {Ñ}{{\~N}}1
}

\newcommand{\captura}[1]{%
  \begin{center}
    \includegraphics[width=\linewidth,height=0.5\textheight,keepaspectratio]{#1}
  \end{center}}

\begin{document}

\begin{titlepage}
\centering

\includegraphics[width=5cm]{logo_usm.png}\\[0.8cm]

{\Large\bfseries Universidad Técnica Federico Santa María}\\[0.2cm]
{\large Departamento de Industrias}

\vspace{1.6cm}

{\Huge\bfseries Laboratorio 4}\\[0.5cm]
{\LARGE\bfseries Bases de Datos Relacionales con MySQL}\\[0.7cm]
{\large Sistemas de Información para la Gestión}

\vspace{1.6cm}

\begin{tabular}{rl}
\textbf{Asignatura:}  & ICN292 Sistemas de Información para la Gestión \\[0.35cm]
\textbf{Profesor:}    & José González \\[0.35cm]
\textbf{Estudiante:}  & Vicente Acevedo \\[0.35cm]
\textbf{Rut:}         & 21.430.418-K \\[0.35cm]
\textbf{Repositorio:} & \url{https://github.com/vacevedo3110/ICN292-Lab4-Acevedo-Vicente} \\
\end{tabular}

\vfill

{\large 7 de noviembre de 2026}

\end{titlepage}


\section*{Consulta 1}

\begin{lstlisting}
Select nombre, genero
FROM autores
where pais <> 'Chile' or pais IS NULL
ORDER by nombre;
\end{lstlisting}

\captura{CONSULTA 1.png}

\noindent La comparación \texttt{pais <>~\textquotesingle Chile\textquotesingle} devuelve verdadero para los países distintos de Chile y desconocido (NULL) para Kenji Arakawa, por eso se agrega \texttt{OR pais IS NULL} para que contabilice todos los paises que no son Chile. 

\newpage
\section*{Consulta 2}

\begin{lstlisting}
SELECT titulo, paginas
from libros
where paginas between 190 and 220
order by paginas;
\end{lstlisting}

\captura{Consulta 2.png}

\newpage
\section*{Consulta 3}

\begin{lstlisting}
Select count(*) AS total_libros,
SUM(ejemplares_vendidos) AS total_ejemplares,
AVG(paginas) AS promedio_paginas,
min(anio) AS anio_mas_antiguo,
max(ejemplares_vendidos) AS max_ejemplares_vendidos
FROM libros;
\end{lstlisting}

\captura{Consulta 3 .png}

\newpage
\section*{Consulta 4}

\begin{lstlisting}
Select  Distinct genero
from autores
order by genero;
\end{lstlisting}

\captura{Consulta 4.png}

\newpage
\section*{Consulta 5}

\begin{lstlisting}
SELECT a.nombre, COUNT(*) AS libros_2024
FROM   autores a
JOIN   libros l ON a.id_autor = l.id_autor
WHERE  l.anio = 2024
GROUP BY a.id_autor, a.nombre
HAVING COUNT(*) > 1
ORDER BY a.nombre;
\end{lstlisting}

\captura{Consulta 5 .png}

\newpage
\section*{Consulta 6}

\begin{lstlisting}
Select titulo, ejemplares_vendidos
from libros
where anio =2024 and ejemplares_vendidos> 2000
order by ejemplares_vendidos DESC limit 3;
\end{lstlisting}

\captura{Consulta 6.png}

\newpage
\section*{Consulta 7}

\begin{lstlisting}
Insert into libros (id_autor, titulo, paginas, ejemplares_vendidos, anio)
VALUES( 3, 'Papel Mojado', 260, 50, 2024);

Select * from libros where titulo = 'Papel Mojado';
Delete from libros where titulo = 'Papel Mojado';
Select * from libros where titulo = 'Papel Mojado';
\end{lstlisting}

\captura{Consulta 7.1 .png}

\newpage
\captura{Consulta 7.2.png}

\newpage
\section*{Consulta 8}

\begin{lstlisting}
SELECT titulo, ejemplares_vendidos FROM libros WHERE titulo = 'Cielo de Enero';

UPDATE libros
SET    ejemplares_vendidos = ejemplares_vendidos + 100
WHERE  titulo = 'Cielo de Enero';

SELECT titulo, ejemplares_vendidos FROM libros WHERE titulo = 'Cielo de Enero';
\end{lstlisting}

\captura{Consulta 8.png}

\newpage
\captura{Consulta 8.2.png}
\captura{Consulta 8.3.png}

\newpage
\section*{Consulta 9}

\begin{lstlisting}
INSERT INTO colecciones (titulo, fecha_publicacion) VALUES
('Narrativa Austral',   '2022-03-15'),
('Pensamiento y Verso', '2022-09-01'),
('Austral Infantil',    '2023-04-23');

UPDATE libros SET id_coleccion = 1 WHERE id_autor = 1 OR id_autor = 5;
UPDATE libros SET id_coleccion = 2 WHERE id_autor = 2 OR id_autor = 3;
UPDATE libros SET id_coleccion = 3 WHERE id_autor = 4;

SELECT l.titulo            AS libro,
       c.titulo            AS coleccion,
       c.fecha_publicacion AS fecha_coleccion
FROM   libros l
JOIN   colecciones c ON l.id_coleccion = c.id_coleccion
ORDER BY c.fecha_publicacion, l.titulo;
\end{lstlisting}

\captura{Consulta 9 .png}

\newpage
\section*{Consulta 10}

\begin{lstlisting}
INSERT INTO pagos (id_autor, fecha_pago, monto) VALUES
(4, '2026-03-31', 850000.00),
(1, '2026-04-15', 420500.50),
(4, '2026-06-30', 610250.75),
(5, '2026-07-10', 735000.00);

-- 10a: cada pago con autor, fecha y monto
SELECT a.nombre AS autor, p.fecha_pago, p.monto
FROM   pagos p
JOIN   autores a ON p.id_autor = a.id_autor
ORDER BY p.fecha_pago;
\end{lstlisting}

\captura{Consulta 10a.png}

\newpage
\begin{lstlisting}
-- 10b: total pagado a cada autor
SELECT a.nombre AS autor, SUM(p.monto) AS total_pagado
FROM   pagos p
JOIN   autores a ON p.id_autor = a.id_autor
GROUP BY a.id_autor, a.nombre
ORDER BY total_pagado DESC;
\end{lstlisting}

\captura{Consulta 10b.png}

\end{document}